\documentclass[12pt, titlepage]{article}
\usepackage{hyperref}

\title{Feedback Controls Lab Weekly Report \\ \normalsize{Week 3}}
\author{Aydan Vissing, Brady Schick, Gabriel Southwell, Simon Ross}
\date{\today}

\begin{document}
\maketitle

\section*{Aydan Vissing (8 hrs spent)}

Spent some time on Wednesday researching some more drone kits (this was before we knew that we weren’t going to be buying one). Thursday we did the PCB lab and Friday we had a meeting to discuss the further development of the proposal and the drones. Saturday we spent a lot of time working on the proposal. I was assigned to finish the Gantt Chart, which I should have done this Wednesday. 

\section*{Brady Schick (8 hrs spent)}

On Thursday, we completed the PCB lab. On Friday, we met to discuss the project proposal. On Saturday, we spent a lot of time putting together our project proposal. I also spent a lot of time doing some independent research.

\section*{Gabriel Southwell (9.5 hrs spent)}

This week I tentatively settled on using the ESP32-S3 for the drone’s MPU. Previously I was concerned that the PWM pins were not going to predictable enough for drone motors. That is why was looking at an STM board since they often have extra drone hardware. This week I came across a few open source drone platforms that worked fine being controlled just by an ESP. This does reduce the overall complexity and some of the cost a bit.

There are many different versions of ESP32, but I think the S3 is a good starting place. It has 2 cores, good networking, and an FPU. It can also be had in multiple different RAM and EEPROM capacities. For prototyping I think a mid range board will be good enough. When its done we might be able to put cheaper ones in the kits. Mid-range boards run about \$12-15.

My team met for around 4 hours on Saturday to work on the project proposal and migrate our workspace and stack to git.
	
On Thursday I finished the PCB MP, and today I worked on my team’s book lecture for Wednesday.

\section*{Simon Ross (9 hrs spent)}

This week was a particularly front-loaded week as homework goes, but at the end of the week, we were able to accomplish a significant amount towards the end goals of our project. It has been interesting to balance working on the project with writing a proposal simultaneously, but I think our team has been successful at staying on the same page and making wise decisions. Additionally, we were able to accomplish a bunch of work in related areas as well.

My main work this week consisted of completing the PCB workshop with Dr. Kohl, brainstorming and organizing our proposal and plan, and working on our presentation for Senior Design class on Wednesday. The PCB workshop was simple and quick, but will likely be very useful later in our project. On Saturday, our team met to create a basic plan that considered financial, temporal, software, and hardware needs for the entire project. This work is directly related to and provides an excellent foundation for the project proposal that we will submit at the end of this week. We also transitioned our workspace from Google Drive to the far more versatile GitHub. This allows us to work in tandem on \LaTeX documents and code that Google is incompatible with. Finally, we switched presentation groups with the NASA team, which has made this week super busy. 

In summary, the project is progressing quite well. We have tentatively decided on a microcontroller, made excellent progress on research, and completed much of the necessary thought for our project proposal (and thus planning in general). This upcoming week will include much more work on our proposal, potentially finishing up with ordering some physical components. And of course, we will present our reading from the Design textbook.

\end{document}